\section{CONCLUSION}

In this paper, we characterized the conditions under which a predictive model benefits from permitting voluntary feature disclosure rather than enforcing mandatory disclosure. Focusing on the regression setting, we framed the strategic interaction as a sender-receiver game and established formal sufficiency and necessity conditions for the existence of a strategic advantage. These conditions highlight the critical role of preference misalignment and the impact of the noisy transmission channel parameter $q$ on the receiver's optimal disclosure regime design. 

\subsection{Limitations}
Our current analysis operates under certain assumptions that limit its generality. First, the derived sufficiency bounds, while theoretically sound, can be conservative and may not tightly capture the exact transition point where voluntary disclosure becomes beneficial in all distribution families. Second, the theoretical framework currently restricts the receiver to linear strategies. While this choice provides tractability and clean, interpretable algebraic conditions, it restricts the receiver from deploying highly expressive, non-linear predictors that could potentially extract more nuanced information from opportunistic reporting. 

\subsection{Future Work}
There are several promising directions for extending this research. A natural next step is to adapt these bounds beyond continuous regression to handle discrete problem formulations, such as binary classification, where the interplay of conditional distributions may lead to different threshold behaviors. Furthermore, our analysis assumes the environment and distributions are perfectly known to the receiver. Extending our framework to an online learning or dynamic setting---where the receiver must interact repeatedly with senders to learn the environment parameters while minimizing cumulative regret---remains a compelling direction. Finally, extending the game-theoretic formulation to accommodate multi-sender environments, where strategic feature disclosure from one entity influences the competitive dynamics of others, represents a challenging but highly applicable frontier.
